\subsection*{Принцип ДП по подотрезкам}

Если имеется способ расчета ответ на задачу на основе ответов на подотреках, то выгодно применить ДП по подотрезкам.

То есть для решения таковой задачи надо сделать следующее.

\begin{enumerate}
    \item $dp[l,\,r]$ — ответ на задачу на подотрезке $[l,\,r]$
    \item Задать базу для отрезков длины 1, 2
    \item Определить функцию «слияния» ответов. То есть как получить из результата на подотрезках результат для всего отрезка.
\end{enumerate}

\subsection*{Квадратичное ДП по подотрезкам}

Данная постановка предполагает, что ответ для отрезка $[l,\,r]$ можно извлечь из отрезков $[l + 1,\,r]$, $[l,\,r - 1]$, $[l + 1,\,r - 1]$.

\subsubsection*{Пример задачи}

Есть строка $S$. Надо найти длину максимальной подпоследовательности палиндрома за $O(|S|^2)$.

\Solve

\begin{enumerate}
    \item \textit{Смысл dp:} $dp[l][r]$ -- ответ для строки $S[l:r]$.
    \item \textit{База:} $dp[l][l] = 1$ и $dp[l][r] = 0$, если $r < l$.
    \item \textit{Ф-ла пересчета:} У нас могут быть две ситуации:
    \begin{itemize}
        \item Если $S[l] = S[r]$, то тогда нам выгодно оба эти символа докинуть в ответ и $dp[l][r]\,=\,dp[l + 1][r - 1]\,+\,2$.
        \item Если $S[l] \neq S[r]$, то тогда нам нет смысла учитывать данную букву, так что $dp[l][r]\,=\,max \{ dp[l + 1][r],\,dp[l][r - 1] \}$
    \end{itemize}
    \item \textit{Порядок вычислений:} $for\,(len=2 \dots n)\, \{ for\,(l = 0 \dots n) \}$
    \item \textit{Ответ:} $dp[0][|S| - 1]$
\end{enumerate}

\begin{lstlisting}
for (int len = 2; len < S.size(); ++len) {
  for (int l = 0; l < S.size() - len; ++l) {
    int r = l + len;
    if (S[l] == S[r]) {
      dp[l][r] = dp[l + 1][r - 1] + 2;
    } else {
      dp[l][r] = max(dp[l][r - 1], dp[l + 1][r]);
    }
  }
}
\end{lstlisting}

\subsection*{Кубическое ДП по подотрезкам}
Данная постановка предполагает, что ответ для отрезка $[l,\,r]$ можно извлечь из отрезков $[l,\,m]$, $[m + 1,\,r]$, где $m \in [l,\,r]$.

\subsubsection*{Пример задачи}
Есть строка S из разных типов скобок. Надо найти длину максимальной правильной скобочной подпоследовательности (ПСП) за $O(|S|^3)$.

\Solve

\begin{enumerate}
    \item \textit{Смысл dp:} $dp[l][r]$ -- минимальное число символов, которые надо удалить из $S[l:r + 1]$, чтобы получилась ПСП.
    \item \textit{База:} $dp[l][l] = 1$ и $dp[l][r] = 0$, если $r < l$.
    \item \textit{Ф-ла пересчета:} Посмотрим, что может произойти с $S[l]$
    \begin{itemize}
        \item Скобка будет удалена, тогда $dp[l][r]\,=\,dp[l+1][r]\,+\,1$.
        \item Скобка не будет удалена, тогда найдется парная ей скобка на
позиции $m \in [l + 1,\,r - 1]$.\\
Тогда итоговая ПСП на $S[l:r + 1]$ будет устроена как $(ans[l + 1 : m])ans[m + 1 : r + 1]$, то есть $dp[l][r]\,=\,dp[l + 1][m - 1]\,+\,dp[m + 1][r]$.\\
Тогда $dp[l][r]\,=\, \min\limits_{m \in (l,\,r)} ( dp[l + 1][m - 1]\,+\,dp[m + 1][r]$ ).
    \end{itemize}
    \item \textit{Порядок вычислений:} $for\,(len=2 \dots n)\, \{ for\,(l = 0 \dots n) \}$
    \item \textit{Ответ:} $|S|\,-\,dp[0][|S| - 1]$.
\end{enumerate}

\begin{lstlisting}
for (int len = 2; len < S.size(); ++len) {
  for (int l = 0; l < S.size() - len; ++l) {
    int r = l + len;
    dp[l][r] = dp[l + 1][r] + 1;
    for (int m = l + 1; m < r; ++m) {
      if (S[l] == '(' && S[m] == ')') {
        dp[l][r] = min(dp[l][r], dp[l + 1][m - 1] + dp[m + 1][r]);
      }
    }
  }
}
\end{lstlisting}

\subsection*{Задача об оптимальном перемножении матриц}

\Task Дана последовательность из $n$ матриц, требуется найти самый эффективный способ их перемножения.

У нас есть множество способов перемножить матрицы, потому что операция перемножения ассоциативна. Другими словами, нет разницы в каком порядке расставляются скобки между множителями, результат будет один и тот же. Однако, порядок в котором расставляются скобки между матрицами повлияет на количество арифметических операций, которые потребуются на вычисление ответа

Например, предположим, что $\dim{A}= 10 \times 30$, $\dim{B} = 30 \times 5$, $\dim{C} = 5 \times 60$. Тогда:\\
\tabДля $(A \times B)\times C$ будет $(10\times30\times5) + (10\times5\times60)  = 1500 + 3000 = 4500$ операций \\
\tabДля $ A \times(B \times C)$ будет $(30\times5\times60) + (10\times30\times60) = 9000 + 18000 = 27000$ операций.

Как мы видим, первый способ гораздо эффективней. 

\Solve 

Пусть $p_i$ -- размерность матрицы $i$

\begin{enumerate}
    \item \textit{Смысл dp:} $dp[l][r]$ -- минимальная стоимость вычисления матрицы $M_{l \dots r}$
    \item \textit{База:} $dp[l][r] = 
    \begin{cases}
    0, & l = r\\
    p_{l-1} \cdot p_{l} \cdot p_{l+1}, & r = l + 1\\
    \end{cases}$.
    \item \textit{Ф-ла пересчета:} $dp[l][r] = 
        \begin{cases}
        0, & l = r\\
        \min\limits_{m \in (l,\,r)} ( dp[l][m] + dp[m + 1][r] + p_{l-1} \cdot p_m \cdot p_r), & r < l
        \end{cases}$
    \item \textit{Порядок вычислений:} $for\,(len=2 \dots n)\, \{ for\,(l = 0 \dots n) \}$
    \item \textit{Ответ:} $dp[0][n-1]$
\end{enumerate}

\begin{lstlisting}
for (int i = 1; i < n; ++i) {
  dp[i][i] = 0;
}
for (int len = 2; len < n; ++len) {
  for (int l = 0; i < n - len; ++l) {
    int r = i + len;
    dp[l][r] = +inf;
    for (int m = l + 1; m < r; ++m) {
      dp[l][r] = min(dp[l][r], dp[l][m] + dp[m + 1][r] + p[l - 1] * p[m] * p[r]);
    }
  }
}
\end{lstlisting}

\pagebreak